\documentclass[aspectratio=169,10pt]{beamer}
% Shared preamble for the Algorithms and Programming Concepts slides (Beamer, Metropolis)
\usetheme[progressbar=frametitle,numbering=fraction,background=light,block=fill,sectionpage=none]{metropolis}
\usepackage[T1]{fontenc}
\usepackage{FiraSans}
\usepackage[scaled=.85]{FiraMono}
\usepackage{amsmath,amssymb}
\usefonttheme[onlymath]{serif}
\usepackage{booktabs,array,tabularx}
\usepackage{tikz,pgfplots}
\usepackage{xcolor}
\usepackage{listings}
\usepackage{adjustbox}
\usepackage{fontawesome5}
\hyphenpenalty=10000\exhyphenpenalty=10000
\newenvironment{sfig}[1][0.80]{\begin{adjustbox}{max width=\linewidth,max totalheight=#1\textheight,center}}{\end{adjustbox}}
\usetikzlibrary{arrows.meta,positioning,shapes.geometric,shapes.misc,fit,calc,backgrounds,decorations.pathreplacing,patterns,matrix,chains}
\pgfplotsset{compat=1.17}

\definecolor{tdInk}{HTML}{1D2330}
\definecolor{tdMuted}{HTML}{5D6678}
\definecolor{tdLine}{HTML}{D9D5CA}
\definecolor{tdPaper}{HTML}{FBFAF7}
\definecolor{tdBlue}{HTML}{2457D6}
\definecolor{tdGreen}{HTML}{17855A}
\definecolor{tdAmber}{HTML}{B7791F}
\definecolor{tdRed}{HTML}{C8362F}
\definecolor{tdPurple}{HTML}{7A4FD1}
\setbeamercolor{normal text}{fg=tdInk,bg=tdPaper}
\setbeamercolor{alerted text}{fg=tdRed}
\setbeamercolor{example text}{fg=tdGreen}
\setbeamercolor{frametitle}{bg=tdInk,fg=white}
\setbeamercolor{progress bar}{fg=tdBlue,bg=tdLine}
\setbeamercolor{title separator}{fg=tdBlue}
\setbeamercolor{block title}{bg=tdBlue!12,fg=tdInk}
\setbeamercolor{block body}{bg=tdBlue!5,fg=tdInk}
\setbeamercolor{block title alerted}{bg=tdRed!12,fg=tdRed!80!black}
\setbeamercolor{block body alerted}{bg=tdRed!5,fg=tdInk}
\setbeamercolor{block title example}{bg=tdGreen!14,fg=tdGreen!70!black}
\setbeamercolor{block body example}{bg=tdGreen!6,fg=tdInk}
\setbeamerfont{frametitle}{size=\large,series=\bfseries}
\makeatletter\setlength{\metropolis@frametitle@padding}{1.15ex}\makeatother
\setbeamertemplate{navigation symbols}{}
\setbeamerfont{block title}{size=\small}

% ---- code listings (Python 3.14) ----
\lstdefinestyle{py}{language=Python,
  basicstyle=\ttfamily\fontsize{8.4}{10.1}\selectfont\color{tdInk},
  keywordstyle=\color{tdBlue}\bfseries, commentstyle=\color{tdMuted}\itshape, stringstyle=\color{tdGreen},
  morekeywords={match,case,type,self,True,False,None},
  emph={print,len,range,sorted,min,max,sum,enumerate,zip,heappush,heappop,cache},emphstyle=\color{tdPurple},
  showstringspaces=false, columns=fullflexible, keepspaces=true, frame=single, rulecolor=\color{tdLine},
  backgroundcolor=\color{white}, xleftmargin=2pt, xrightmargin=2pt, aboveskip=2pt, belowskip=0pt, framesep=3pt}
\lstdefinestyle{out}{basicstyle=\ttfamily\fontsize{8}{9.6}\selectfont\color{tdInk}, columns=fullflexible, keepspaces=true,
  frame=single, rulecolor=\color{tdGreen!60}, backgroundcolor=\color{tdGreen!6}, xleftmargin=2pt, xrightmargin=2pt, aboveskip=2pt, belowskip=0pt, framesep=3pt}
% \pycode[size]{snippet} shows snippets/<snippet>.py ; \pyout[size]{snippet} shows snippets/<snippet>.out
\newcommand{\pycode}[2][\relax]{{#1\lstinputlisting[style=py]{snippets/#2.py}}}
\newcommand{\pyout}[2][\relax]{{#1\lstinputlisting[style=out]{snippets/#2.out}}}

% ---- tikz styles ----
\tikzset{>={Stealth[length=2.2mm]},
  card/.style={rounded corners=4pt,draw=#1!70!black,fill=#1!10,align=center,inner sep=5pt,font=\small}, card/.default=tdBlue,
  pill/.style={rounded corners=7pt,draw=none,fill=#1,text=white,inner sep=3pt,font=\scriptsize\bfseries,align=center}, pill/.default=tdBlue,
  arr/.style={->,thick,draw=tdMuted},
  cell/.style={draw=tdInk!70,minimum width=7.5mm,minimum height=7.5mm,inner sep=0pt,font=\small\ttfamily,fill=white},
  tn/.style={circle,draw=tdInk!80,fill=white,minimum size=7mm,inner sep=0pt,font=\small\bfseries}}
\newcommand{\key}[1]{\textbf{\color{tdBlue}#1}}
\newcommand{\bad}[1]{\textbf{\color{tdRed}#1}}
\newcommand{\good}[1]{\textbf{\color{tdGreen!70!black}#1}}
\newcommand{\apcfooter}[1]{\setbeamertemplate{footline}{\hfill\tiny\color{tdMuted}Algorithms and Programming Concepts \,|\, #1 \,|\, BUV\hspace{1.6em}\insertframenumber/\inserttotalframenumber\hspace{0.8em}\vspace{4pt}}}
\newcommand{\apctitle}[2]{\title{#1}\subtitle{#2}\author{Dr.\ Quang-Vinh Dang\\ {\small Senior Lecturer (Assoc.\ Prof)}}\institute{British University Vietnam}\date{Algorithms and Programming Concepts \,|\, Year 1 \,|\, Python 3.14}}

\apcfooter{Graphs}
\apctitle{Graphs}{Modelling connections: BFS, DFS and shortest paths}

\tikzset{gn/.style={rounded corners=3pt,draw=tdInk!70,minimum width=11mm,minimum height=6.5mm,inner sep=2pt,font=\scriptsize\bfseries,fill=white},
         wl/.style={midway,fill=tdPaper,inner sep=1pt,font=\tiny\bfseries,text=tdMuted}}
\newcommand{\gA}{white}\newcommand{\gB}{white}\newcommand{\gC}{white}\newcommand{\gD}{white}\newcommand{\gE}{white}\newcommand{\gF}{white}
% campus graph; #1 = 1 shows the walking minutes on the edges
\newcommand{\campus}[1]{%
  \coordinate (pg) at (0,0); \coordinate (pl) at (2.6,1.5); \coordinate (pc) at (2.6,-1.5);
  \coordinate (pb) at (5.6,1.5); \coordinate (pd) at (5.6,-1.5); \coordinate (py) at (8.2,0);
  \ifnum#1=1
    \draw[tdInk!60] (pg) -- node[wl] {4} (pl); \draw[tdInk!60] (pg) -- node[wl] {2} (pc);
    \draw[tdInk!60] (pl) -- node[wl] {5} (pb); \draw[tdInk!60] (pl) -- node[wl,pos=0.3] {6} (pd);
    \draw[tdInk!60] (pc) -- node[wl] {3} (pd); \draw[tdInk!60] (pb) -- node[wl] {6} (py); \draw[tdInk!60] (pd) -- node[wl] {2} (py);
  \else
    \foreach \u/\v in {pg/pl,pg/pc,pl/pb,pl/pd,pc/pd,pb/py,pd/py}{\draw[tdInk!60] (\u) -- (\v);}
  \fi
  \node[gn,fill=\gA] (gate) at (pg) {Gate}; \node[gn,fill=\gB] (lib) at (pl) {Lib}; \node[gn,fill=\gC] (cafe) at (pc) {Cafe};
  \node[gn,fill=\gD] (lab) at (pb) {Lab}; \node[gn,fill=\gE] (dorm) at (pd) {Dorm}; \node[gn,fill=\gF] (gym) at (py) {Gym};}
\newcommand{\badge}[3]{\node[pill=#3,inner sep=2pt,font=\tiny\bfseries] at ($(#1)+(0.5,0.42)$) {#2};}

\begin{document}
\begin{frame}[plain]\titlepage\end{frame}

\begin{frame}{Graphs model connections}
\begin{columns}[T]
\begin{column}{0.5\textwidth}
{\small
A \key{graph} is a set of things (\textbf{nodes}) and the links between them (\textbf{edges}). It is the most flexible structure in computing:
\begin{itemize}\setlength\itemsep{3pt}
\item \textbf{maps}: places and roads (Google Maps);
\item \textbf{social networks}: people and friendships;
\item \textbf{the web}: pages and hyperlinks;
\item \textbf{dependencies}: which course needs which;
\item \textbf{networks}: routers and cables.
\end{itemize}
Lists and trees are special cases of graphs.}
\end{column}
\begin{column}{0.46\textwidth}
\begin{block}{After this lecture you can}
{\small
\begin{itemize}\setlength\itemsep{2pt}
\item represent a graph with a dictionary (adjacency list) or a matrix;
\item explore it with \textbf{BFS} and \textbf{DFS};
\item find the \textbf{fewest-hops} route and the \textbf{cheapest} route (Dijkstra);
\item find \textbf{connected components};
\item say what each costs: $O(V+E)$, $O((V+E)\log V)$.
\end{itemize}}
\end{block}
\end{column}
\end{columns}
\end{frame}

\begin{frame}{Our running example: a small campus}
\begin{columns}[T]
\begin{column}{0.6\textwidth}
\begin{sfig}[0.6]
\begin{tikzpicture}
  \campus{0}
\end{tikzpicture}
\end{sfig}
\end{column}
\begin{column}{0.38\textwidth}
{\small
\begin{description}\setlength\itemsep{2pt}
\item[vertex / node] a place: Gate, Lib, \ldots
\item[edge] a footpath between two places.
\item[neighbours] nodes joined by an edge: Lib's neighbours are Gate, Lab, Dorm.
\item[degree] number of neighbours: Dorm has 3.
\item[path] a walk along edges.
\item[cycle] a path back to its start.
\item[connected] every node can reach every other.
\end{description}}
\end{column}
\end{columns}
\vspace{1mm}
{\small This graph is \textbf{undirected} (paths work both ways) and, later, \textbf{weighted} (each path takes some minutes). One-way streets would make it \textbf{directed}.}
\end{frame}

\begin{frame}{Two ways to store a graph}
\begin{columns}[T]
\begin{column}{0.5\textwidth}
\pycode{representation}
\vspace{2mm}
\pyout{representation_demo}
\vspace{1mm}
{\small \textbf{Adjacency list}: a dict from each node to the list of its neighbours.}
\end{column}
\begin{column}{0.48\textwidth}
\pyout{matrix_demo}
\vspace{1mm}
{\small \textbf{Adjacency matrix}: a table with 1 where an edge exists (built by the code below).}
\vspace{1mm}
\pycode{matrix}
\end{column}
\end{columns}
\end{frame}

\begin{frame}{Which representation?}
\renewcommand{\arraystretch}{1.35}\small
\begin{tabular}{@{}>{\raggedright\bfseries}p{50mm}>{\raggedright}p{45mm}>{\raggedright\arraybackslash}p{45mm}@{}}
\toprule
 & Adjacency list & Adjacency matrix\\\midrule
Memory & $O(V+E)$ & $O(V^2)$\\
Is there an edge $u\!-\!v$? & scan $u$'s list: $O(\text{degree})$ & one lookup: $O(1)$\\
List all neighbours of $u$ & $O(\text{degree})$ & scan a row: $O(V)$\\
Best for & \textbf{sparse} graphs (most real ones) & \textbf{dense} graphs, quick edge tests\\
\bottomrule
\end{tabular}
\vspace{4mm}

{\small A social network with a billion people where each knows a few hundred others has $E\approx 10^{11}$ but $V^2=10^{18}$ cells. The matrix would not fit in any computer, so real systems use adjacency lists. We use them from now on.}
\end{frame}

\begin{frame}{Breadth-first search (BFS): explore in layers}
\begin{columns}[T]
\begin{column}{0.58\textwidth}
\begin{sfig}[0.55]
\begin{tikzpicture}
  \renewcommand{\gA}{tdAmber!40}\renewcommand{\gB}{tdBlue!22}\renewcommand{\gC}{tdBlue!22}\renewcommand{\gD}{tdGreen!24}\renewcommand{\gE}{tdGreen!24}\renewcommand{\gF}{tdPurple!22}
  \campus{0}
  \node[font=\tiny,text=tdAmber!70!black,anchor=south] at (0,0.5) {0 hops};
  \node[font=\tiny,text=tdBlue,anchor=south] at (2.6,2.0) {1 hop}; \node[font=\tiny,text=tdBlue,anchor=north] at (2.6,-2.0) {1 hop};
  \node[font=\tiny,text=tdGreen!70!black,anchor=south] at (5.6,2.0) {2 hops}; \node[font=\tiny,text=tdGreen!70!black,anchor=north] at (5.6,-2.0) {2 hops};
  \node[font=\tiny,text=tdPurple,anchor=south] at (8.2,0.5) {3 hops};
\end{tikzpicture}
\end{sfig}
\end{column}
\begin{column}{0.4\textwidth}
{\small
Start at the Gate. Visit everything \textbf{1 hop} away, then everything \textbf{2 hops} away, and so on: the search spreads like a ripple.

\vspace{2mm}
To do this we keep a \textbf{queue} (first in, first out): neighbours are added at the back, and we always take the node at the front.

\vspace{2mm}
\textbf{Bonus:} the first time BFS reaches a node, it has used the \textbf{fewest possible edges}. So BFS finds shortest paths in an \emph{unweighted} graph.}
\end{column}
\end{columns}
\end{frame}

\begin{frame}{BFS in Python}
\begin{columns}[T]
\begin{column}{0.5\textwidth}
\pycode{bfs}
\vspace{2mm}
{\small
\texttt{parent} does two jobs: it remembers which nodes we have already seen (so none is queued twice), and it records \emph{how we got there}.}
\end{column}
\begin{column}{0.47\textwidth}
\pycode{path}
\vspace{2mm}
\pyout{bfs_demo}
\vspace{2mm}
{\small
Following \texttt{parent} backwards from Gym gives the route. Every node enters the queue once and every edge is looked at once: $O(V+E)$.}
\end{column}
\end{columns}
\end{frame}

\begin{frame}{Depth-first search (DFS): go deep, then backtrack}
\begin{columns}[T]
\begin{column}{0.5\textwidth}
\begin{sfig}[0.5]
\begin{tikzpicture}
  \campus{0}
  \badge{pg}{1}{tdAmber}\badge{pl}{2}{tdAmber}\badge{pb}{3}{tdAmber}\badge{py}{4}{tdAmber}\badge{pd}{5}{tdAmber}\badge{pc}{6}{tdAmber}
  \draw[->,very thick,tdAmber!80!black] (pg) -- (pl); \draw[->,very thick,tdAmber!80!black] (pl) -- (pb);
  \draw[->,very thick,tdAmber!80!black] (pb) -- (py); \draw[->,very thick,tdAmber!80!black] (py) -- (pd); \draw[->,very thick,tdAmber!80!black] (pd) -- (pc);
\end{tikzpicture}
\end{sfig}
{\small Follow one path as far as it goes; when you hit a dead end (or a place already visited), \textbf{backtrack} to the last choice and try the next neighbour.}
\end{column}
\begin{column}{0.48\textwidth}
\pycode{dfs}
\vspace{2mm}
\pyout{dfs_demo}
\vspace{2mm}
{\small
Recursion \emph{is} the stack: each call remembers where to continue. Also $O(V+E)$. DFS does not find shortest paths, but it is ideal for ``can I get there?'', cycles, mazes and ordering tasks.}
\end{column}
\end{columns}
\end{frame}

\begin{frame}{Connected components: separate islands}
\begin{columns}[T]
\begin{column}{0.5\textwidth}
\pycode{components}
\vspace{2mm}
\pyout{components_demo}
\end{column}
\begin{column}{0.46\textwidth}
\begin{sfig}[0.4]
\begin{tikzpicture}
  \node[gn,fill=tdBlue!22] (a) at (0,0.5) {A}; \node[gn,fill=tdBlue!22] (b) at (1.6,0.5) {B}; \node[gn,fill=tdBlue!22] (c) at (3.2,0.5) {C};
  \node[gn,fill=tdGreen!24] (d) at (0,-0.9) {D}; \node[gn,fill=tdGreen!24] (e) at (1.6,-0.9) {E}; \node[gn,fill=tdAmber!35] (f) at (3.2,-0.9) {F};
  \draw[tdInk!60] (a) -- (b) -- (c); \draw[tdInk!60] (d) -- (e);
\end{tikzpicture}
\end{sfig}
{\small
A \textbf{component} is a group of nodes that can all reach each other. Start a DFS from any unseen node: everything it reaches is one component; repeat until all nodes are seen.

\vspace{2mm}
Uses: ``are these two accounts connected?'', finding islands in an image, checking that a network is not split in two.}
\end{column}
\end{columns}
\end{frame}

\begin{frame}{Fewest hops is not the same as quickest}
\begin{columns}[T]
\begin{column}{0.54\textwidth}
\begin{sfig}[0.52]
\begin{tikzpicture}
  \campus{1}
  \foreach \u/\v in {pg/pl,pl/pb,pb/py}{\draw[very thick,tdRed!80,opacity=0.75] (\u) -- (\v);}
  \foreach \u/\v in {pg/pc,pc/pd,pd/py}{\draw[very thick,tdGreen!80!black,opacity=0.75] (\u) -- (\v);}
  \campus{1}
\end{tikzpicture}
\end{sfig}
{\small
Both routes from Gate to Gym have \textbf{3 edges}, so BFS cannot tell them apart. It returned Gate--Lib--Lab--Gym:
\textcolor{tdRed}{\textbf{red}} $4+5+6=\textbf{15}$ min, but \textcolor{tdGreen!70!black}{\textbf{green}} takes $2+3+2=\textbf{7}$ min.}
\end{column}
\begin{column}{0.43\textwidth}
{\small The weighted graph in Python: a dict of dicts (neighbour $\to$ minutes to walk).}
\vspace{1mm}
\pycode[\fontsize{7.8}{9.4}\selectfont]{weighted}
\vspace{2mm}
{\small
When edges have different costs we need an algorithm that adds up the \emph{weights}: \key{Dijkstra's algorithm}.}
\end{column}
\end{columns}
\end{frame}

\begin{frame}{Dijkstra's algorithm: settle the closest node first}
\begin{columns}[T]
\begin{column}{0.46\textwidth}
\begin{sfig}[0.5]
\begin{tikzpicture}
  \renewcommand{\gA}{tdAmber!40}\renewcommand{\gB}{tdBlue!15}\renewcommand{\gC}{tdBlue!15}\renewcommand{\gD}{tdBlue!15}\renewcommand{\gE}{tdBlue!15}\renewcommand{\gF}{tdBlue!15}
  \campus{1}
  \node[font=\tiny\bfseries,text=tdGreen!70!black,anchor=south] at (0,0.45) {0};
  \node[font=\tiny\bfseries,text=tdGreen!70!black,anchor=south] at (2.6,1.95) {4};  \node[font=\tiny\bfseries,text=tdGreen!70!black,anchor=north] at (2.6,-1.95) {2};
  \node[font=\tiny\bfseries,text=tdGreen!70!black,anchor=south] at (5.6,1.95) {9};  \node[font=\tiny\bfseries,text=tdGreen!70!black,anchor=north] at (5.6,-1.95) {5};
  \node[font=\tiny\bfseries,text=tdGreen!70!black,anchor=south] at (8.2,0.45) {7};
\end{tikzpicture}
\end{sfig}
{\small Green numbers = final shortest time from Gate.}
\vspace{2mm}

\pyout{dijkstra_trace}
\end{column}
\begin{column}{0.5\textwidth}
{\small
\begin{enumerate}\setlength\itemsep{3pt}
\item Give every node a distance: 0 for the start, $\infty$ for the rest.
\item Take the \textbf{unsettled node with the smallest distance}. Its distance is now final.
\item \textbf{Relax} its edges: if going through it makes a neighbour closer, update that neighbour.
\item Repeat until everything is settled.
\end{enumerate}
\vspace{1mm}
A \textbf{heap} (priority queue, last lecture!) gives ``the closest node'' in $O(\log n)$. Why is it correct? All weights are $\ge 0$, so a longer walk can never become shorter.

\vspace{1mm}
\textbf{Negative weights} break this argument: use Bellman--Ford.}
\end{column}
\end{columns}
\end{frame}

\begin{frame}{Dijkstra in Python}
\begin{columns}[T]
\begin{column}{0.5\textwidth}
\pycode{dijkstra}
\end{column}
\begin{column}{0.47\textwidth}
\pyout{dijkstra_demo}
\vspace{2mm}
{\small
\texttt{dist[v]} is the best time found so far to reach \texttt{v}; \texttt{parent} lets us rebuild the route, exactly as in BFS.

\vspace{2mm}
\texttt{if d > dist[u]: continue} skips \emph{stale} heap entries: a node may be pushed several times with different distances, and only the smallest one matters.

\vspace{2mm}
Cost: each edge may push once, each push or pop is $O(\log)$: $O((V+E)\log V)$.}
\end{column}
\end{columns}
\end{frame}

\begin{frame}{Which graph algorithm?}
\renewcommand{\arraystretch}{1.32}\footnotesize
\begin{tabular}{@{}>{\raggedright\bfseries}p{26mm}>{\raggedright}p{50mm}>{\raggedright}p{30mm}>{\raggedright\arraybackslash}p{36mm}@{}}
\toprule
Algorithm & Use it to \ldots & Cost & Data structure\\\midrule
BFS & fewest-edge path, layers, ``how far?'' & $O(V+E)$ & queue (\texttt{deque})\\
DFS & reachability, components, cycles, mazes & $O(V+E)$ & recursion / stack\\
Dijkstra & cheapest path, weights $\ge 0$ & $O((V+E)\log V)$ & heap (\texttt{heapq})\\
Bellman--Ford & cheapest path with negative weights & $O(V\cdot E)$ & arrays\\
A* & cheapest path to \emph{one} target, using a guess of the remaining distance & often much faster than Dijkstra & heap\\
\bottomrule
\end{tabular}
\vspace{4mm}

{\small Navigation apps run Dijkstra-like searches over road graphs with hundreds of millions of edges, with many extra tricks to answer in milliseconds.}
\end{frame}

\begin{frame}{Recap and practice}
\begin{columns}[T]
\begin{column}{0.48\textwidth}
\begin{block}{Take-aways}
{\small
\begin{enumerate}\setlength\itemsep{3pt}
\item A graph is nodes plus edges; store it as an \key{adjacency list} (a dict of lists).
\item \key{BFS} explores in layers with a queue and finds the fewest-hop path.
\item \key{DFS} dives deep with recursion; use it for reachability and components.
\item With edge costs, use \key{Dijkstra}: settle the closest node first, using a heap.
\item BFS and DFS cost $O(V+E)$; Dijkstra $O((V+E)\log V)$.
\end{enumerate}}
\end{block}
\end{column}
\begin{column}{0.48\textwidth}
\begin{block}{Practice}
{\small
\begin{enumerate}\setlength\itemsep{3pt}
\item Add a node \texttt{Pool} connected to \texttt{Gym}. What changes in the BFS output?
\item Write \texttt{has\_path(graph, a, b)} using BFS or DFS.
\item Make Dijkstra return the distance to \emph{one} target and stop early.
\item Does Dijkstra still work if you change one road to 0 minutes? To $-1$? Why?
\item Course prerequisites form a directed graph. Write a function that lists a valid order to take the courses (topological sort).
\end{enumerate}}
\end{block}
\end{column}
\end{columns}
\end{frame}
\end{document}
